% Unit 053 terjemahan Bahasa Indonesia dari chapter4.tex baris 623--722.
% Sumber: Methods of Algebra, Volume 2, commit
% 9a5803ff77dd3257484cb177f851a73770a59dd3 (CC BY 4.0).
% stable-unit-id: o014.aljabr2.chapter4.verdier-localization
% Cakupan: akhir Bagian 4.3, dari konstruksi $S\mathcal{N}$ dan lokalisasi
% Verdier sampai hubungan lokalisasi dengan subkategori; berhenti sebelum
% Bagian 4.4 pada baris 723.

% segment-id: o014.li2.chapter4.verdier-localization.pr001
\begin{proposisi}
	\index[sym1]{SN@$S\mathcal{N}$}
	Misalkan $\mathcal{N}$ subkategori bertriangulasi jenuh dari suatu kategori
	bertriangulasi $\mathcal{D}$ (Definisi--Proposisi~\ref{def:triangulated-subcat},
	Konvensi~\ref{con:saturated-intersection}). Maka
% segment-id: o014.li2.chapter4.verdier-localization.q001
	\[ S\mathcal{N} := \left\{\begin{array}{r|l}
		s \in \Mor(\mathcal{D}) & \exists \;\text{segitiga terbedakan}\; X \xrightarrow{s} Y \to Z \xrightarrow{+1} \\
		& \text{sedemikian sehingga}\; Z \in \Obj(\mathcal{N})
	\end{array}\right\} \]
	merupakan sistem multiplikatif yang kompatibel dengan triangulasi.
\end{proposisi}

% segment-id: o014.li2.chapter4.verdier-localization.pf001
\begin{bukti}
	Karena $\mathcal{N}$ tertutup terhadap pergeseran, dengan merotasi segitiga
	terlihat bahwa $S\mathcal{N}$ memenuhi (ST1). Selanjutnya, tinjau diagram
	komutatif dalam (ST2) dan andaikan $\alpha, \beta \in S\mathcal{N}$.
	Lema~\ref{prop:triangulated-9-diag} memberikan diagram
% segment-id: o014.li2.chapter4.verdier-localization.g001
	\[\begin{tikzcd}
		{} & {} & {} & {} \\
		X'' \arrow[u, "+1"] \arrow[r] & Y'' \arrow[u, "+1"] \arrow[r] & Z'' \arrow[u, "+1"] \arrow[r, "+1"] & {} \\
		X \arrow[r] \arrow[u] & Y \arrow[r] \arrow[u] & Z \arrow[r, "+1"] \arrow[u] & {}\\
		X' \arrow[r] \arrow[u, "\alpha"] & Y' \arrow[r] \arrow[u, "\beta"] & Z' \arrow[r, "+1"] \arrow[u, "\gamma"] & {}
	\end{tikzcd}\]
	sedemikian sehingga setiap baris dan kolom merupakan segitiga terbedakan,
	semua persegi komutatif, dan $(\alpha, \beta, \gamma)$ memberikan suatu
	morfisme di antara segitiga-segitiga. Dengan menerapkan sifat jenuh
	$\mathcal{N}$, fakta bahwa $\alpha, \beta \in S$, dan
	Korolari~\ref{prop:Cone-uniqueness}, terlihat bahwa
	$X'', Y'' \in \Obj(\mathcal{N})$. Definisi subkategori bertriangulasi jenuh
	selanjutnya menyiratkan $Z'' \in \Obj(\mathcal{N})$, sehingga
	$\gamma \in S$. Ini memverifikasi (ST2).

	Berikutnya kita verifikasi aksioma-aksioma sistem multiplikatif
	(Definisi~\ref{def:multiplicative-family}). Karena
	$0 \in \Obj(\mathcal{N})$, aksioma (S1) mengikuti dari (TR1).

	Untuk memverifikasi (S2), misalkan $f, g \in S\mathcal{N}$. Pilih segitiga
	terbedakan $X \xrightarrow{f} Y \to Z' \xrightarrow{+1}$ dan
	$Y \xrightarrow{g} Z \to X' \xrightarrow{+1}$ dengan
	$Z', X' \in \Obj(\mathcal{N})$, lalu gunakan (TR2) untuk memilih segitiga
	terbedakan $X \xrightarrow{gf} Z \to Y' \xrightarrow{+1}$. Penerapan (TR5)
	memberikan segitiga terbedakan $Z' \to Y' \to X' \xrightarrow{+1}$.
	Karena itu, sifat subkategori bertriangulasi jenuh menyiratkan
	$Y' \in \Obj(\mathcal{N})$, yakni $gf \in S\mathcal{N}$.

	Untuk memverifikasi (S3), tinjau morfisme-morfisme
	$X \xrightarrow{s \in S\mathcal{N}} Z \xleftarrow{f} Y$. Pilih segitiga
	terbedakan $X \xrightarrow{s} Z \xrightarrow{h} N \xrightarrow{+1}$ dengan
	$N \in \Obj(\mathcal{N})$, lalu gunakan (TR2) untuk memilih segitiga
	$Y \xrightarrow{hf} N \to U \xrightarrow{+1}$. Terapkan (TR4) pada bagian
	bergaris penuh dari diagram
% segment-id: o014.li2.chapter4.verdier-localization.g002
	\[\begin{tikzcd}
		Y \arrow[r, "hf"] \arrow[d, "f"'] & N \arrow[r] \arrow[equal, d] & U \arrow[r, "+1"] \arrow[dashed, d] & TY \arrow[dashed, d, "Tf"] \\
		Z \arrow[r, "h"'] & N \arrow[r] & TX \arrow[r, "-Ts"'] & TZ
	\end{tikzcd}\]
	untuk memperoleh morfisme-morfisme yang ditunjukkan dengan garis putus-putus
	dan dengan demikian suatu morfisme di antara segitiga-segitiga terbedakan.
	Setelah dirotasi, diperoleh unsur $s': W := T^{-1} U \to Y$ dari
	$S\mathcal{N}$, morfisme yang sesuai $f': W \to X$, beserta diagram komutatif
	yang dipersyaratkan dalam (S3):
% segment-id: o014.li2.chapter4.verdier-localization.g003
	\[\begin{tikzcd}
		X \arrow[d, "{s \in S\mathcal{N}}"'] & W \arrow[d, "{s' \in S\mathcal{N}}"] \arrow[l, "{f'}"'] \\
		Z & Y \arrow[l, "f"] .
	\end{tikzcd}\]

	Untuk memverifikasi (S4), dalam aksioma itu gantikan $(f, g)$ dengan
	$(f-g, 0)$. Cukup dibuktikan pernyataan berikut: untuk suatu
	morfisme $f: X \to Y$, jika terdapat $s: Y \to W$ dengan
	$s \in S\mathcal{N}$ dan $sf = 0$, maka terdapat $t: Z \to X$ dengan
	$t \in S\mathcal{N}$ dan $ft = 0$. Untuk membuktikannya, pilih segitiga
	terbedakan $N \to Y \xrightarrow{s} W \xrightarrow{+1}$ dengan
	$N \in \Obj(\mathcal{N})$. Karena $\Hom(X, \cdot)$ merupakan funktor
	kohomologis dan $sf=0$, terdapat $h \in \Hom(X, N)$ sedemikian sehingga
	$f$ mempunyai faktorisasi $X \xrightarrow{h} N \to Y$. Pilih segitiga
	terbedakan $Z \xrightarrow{t} X \xrightarrow{h} N \xrightarrow{+1}$;
	maka $t \in S\mathcal{N}$. Di sisi lain, $ht=0$
	(Lema~\ref{prop:dt-composite-null}) menyiratkan $ft = 0$. Jadi, (S4)
	terpenuhi.

	Argumen untuk versi dual (S3) dan (S4) persis sama.
\end{bukti}

% segment-id: o014.li2.chapter4.verdier-localization.th001
\begin{teorema}[Lokalisasi Verdier]\label{prop:Verdier-localization}
	\index{lokalisasi!Verdier}
	\index[sym1]{DN@$\mathcal{D}/\mathcal{N}$}
	Misalkan $\mathcal{N}$ subkategori bertriangulasi jenuh dari
	$\mathcal{D}$. Definisikan
	$\mathcal{D}/\mathcal{N} := \mathcal{D}\left[ (S\mathcal{N})^{-1}\right]$
	(yang mungkin merupakan kategori ``besar''), dengan struktur bertriangulasi
	dari Proposisi~\ref{prop:localization-triangulated}. Maka funktor lokalisasi
	$Q: \mathcal{D} \to \mathcal{D}/\mathcal{N}$ mempunyai sifat-sifat berikut.
% segment-id: o014.li2.chapter4.verdier-localization.l001
	\begin{enumerate}[(i)]
% segment-id: o014.li2.chapter4.verdier-localization.i001
		\item Untuk setiap morfisme $f: X \to Y$ dalam $\mathcal{D}$, berlaku
		$Qf = 0$ jika dan hanya jika $f$ dapat difaktorkan sebagai
		$X \to N \to Y$ dengan $N \in \Obj(\mathcal{N})$; khususnya, $Q$
		memetakan $\mathcal{N}$ ke $0$.
% segment-id: o014.li2.chapter4.verdier-localization.i002
		\item Untuk setiap kategori prabertriangulasi $\mathcal{E}$ dan funktor
		bertriangulasi $F: \mathcal{D} \to \mathcal{E}$, jika $F$ memetakan
		$\mathcal{N}$ ke $0$, maka $F$ terfaktorkan secara unik sebagai
		$\mathcal{D} \xrightarrow{Q} \mathcal{D}/\mathcal{N}
		\xrightarrow{\overline{F}} \mathcal{E}$, dengan $\overline{F}$ suatu
		funktor bertriangulasi.
% segment-id: o014.li2.chapter4.verdier-localization.i003
		\item Untuk setiap kategori abelian $\mathcal{A}$ dan funktor kohomologis
		$H: \mathcal{D} \to \mathcal{A}$, jika $H$ memetakan $\mathcal{N}$ ke
		$0$, maka $H$ terfaktorkan secara unik sebagai
		$\mathcal{D} \xrightarrow{Q} \mathcal{D}/\mathcal{N}
		\xrightarrow{\overline{H}} \mathcal{A}$, dengan $\overline{H}$ suatu
		funktor kohomologis.
	\end{enumerate}
\end{teorema}

% segment-id: o014.li2.chapter4.verdier-localization.pf002
\begin{bukti}
	Untuk (i), jika $f$ mempunyai faktorisasi $X \to N \to Y$, maka $Qf = 0$.
	Memang, rotasi $N \to 0 \to TN \xrightarrow{+1}$ dari segitiga terbedakan
	$N \xrightarrow{\identity_N} N \to 0 \xrightarrow{+1}$ menunjukkan bahwa
	$N \to 0$ dipetakan menjadi suatu isomorfisme dalam
	$\mathcal{D}/\mathcal{N}$. Sebaliknya, jika $Qf = 0$, maka terdapat morfisme
	$s: M \to X$ dengan $s \in S\mathcal{N}$ dan $fs = 0$
	(Korolari~\ref{prop:fraction-zero}). Menurut definisi $S\mathcal{N}$,
	terdapat diagram komutatif yang setiap barisnya merupakan segitiga
	terbedakan (bagian bergaris penuh):
% segment-id: o014.li2.chapter4.verdier-localization.g004
	\[\begin{tikzcd}
		M \arrow[r, "s"] \arrow[d] & X \arrow[d, "f"] \arrow[r] & N \arrow[r] \arrow[dashed, d] & TM \arrow[dashed, d] \\
		0 \arrow[r] & Y \arrow[r, "{\identity_Y}"'] & Y \arrow[r] & 0
	\end{tikzcd} \quad (N \in \Obj(\mathcal{N})), \]
	Dengan (TR4), diperoleh morfisme-morfisme yang ditunjukkan dengan garis
	putus-putus sehingga seluruh diagram komutatif. Diagram ini memberikan
	faktorisasi $f$ sebagai $X \to N \to Y$.

	Tinjau situasi dalam (ii). Misalkan $s \in S\mathcal{N}$ dan masukkan $s$
	ke dalam segitiga terbedakan
	$X \xrightarrow{s} Y \to N \xrightarrow{+1}$ dengan
	$N \in \Obj(\mathcal{N})$. Maka
	$FX \xrightarrow{Fs} FY \to 0 \xrightarrow{+1}$ merupakan segitiga
	terbedakan dalam $\mathcal{E}$. Korolari~\ref{prop:triangle-0-isom}
	menunjukkan bahwa $Fs$ merupakan isomorfisme. Sifat universal kemudian
	memberikan faktorisasi unik $H = \overline{F} \circ Q$, dengan
	$\overline{F}$ suatu funktor aditif. Mengingat deskripsi struktur
	bertriangulasi pada $\mathcal{D}/\mathcal{N}$
	(Proposisi~\ref{prop:localization-triangulated}), mudah diverifikasi bahwa
	$\overline{F}$ mempertahankan pergeseran dan segitiga; perinciannya
	diserahkan kepada pembaca.

	Tinjau situasi dalam (iii). Sekali lagi, masukkan $s \in S\mathcal{N}$ ke
	dalam segitiga terbedakan $X \xrightarrow{s} Y \to N \xrightarrow{+1}$.
	Dengan menerapkan $H$, diperoleh barisan eksak
	$\underbracket{HT^{-1} N}_{=0} \to HX \xrightarrow{Hs} HY \to
	\underbracket{HN}_{=0}$, sehingga $Hs$ merupakan isomorfisme. Sifat universal
	kemudian memberikan faktorisasi unik $H = \overline{H} \circ Q$, dengan
	$\overline{H}$ suatu funktor aditif. Untuk memverifikasi bahwa $\overline{H}$
	merupakan funktor kohomologis, ingat bahwa segitiga-segitiga terbedakan dalam
	$\mathcal{D}/\mathcal{N}$ isomorfik dengan citra segitiga-segitiga
	terbedakan dalam $\mathcal{D}$.
\end{bukti}

% segment-id: o014.li2.chapter4.verdier-localization.ex001
\begin{contoh}\label{eg:cohomological-functor-multiplicative}
	\index{kuasi-isomorfisme}
	\index[sym1]{SHSNH@$S_H$, $S\mathcal{N}_H$}
	Funktor kohomologis memberikan suatu kelas penting sistem multiplikatif yang
	kompatibel dengan triangulasi. Misalkan $H: \mathcal{D} \to \mathcal{A}$
	suatu funktor kohomologis. Untuk morfisme $f: X \to Y$ dalam $\mathcal{D}$,
	jika $H(T^n f)$ merupakan isomorfisme dalam $\mathcal{A}$ untuk setiap
	$n \in \Z$, maka $f$ disebut $H$-\emph{kuasi-isomorfisme}. Semua
	$H$-kuasi-isomorfisme membentuk himpunan $S_H \subset \Mor(\mathcal{D})$.
	Himpunan ini jelas memenuhi (ST1), sedangkan (ST2) merupakan akibat dari
	barisan eksak panjang suatu funktor kohomologis
	(Catatan~\ref{rem:cohomological-functor-long}) bersama
	Proposisi~\ref{prop:5-lemma}.

	Di sisi lain, definisikan $\mathcal{N}_H$ seperti dalam
	Contoh~\ref{eg:subcat-N-H}. Ini merupakan subkategori bertriangulasi jenuh,
	sehingga $S\mathcal{N}_H$ dapat didefinisikan. Keduanya mempunyai hubungan
	sederhana berikut.
\end{contoh}

% segment-id: o014.li2.chapter4.verdier-localization.pr002
\begin{proposisi}\label{prop:cohomological-functor-multiplicative}
	Misalkan $H: \mathcal{D} \to \mathcal{A}$ suatu funktor kohomologis. Maka
	$S\mathcal{N}_H = S_H$.
\end{proposisi}

% segment-id: o014.li2.chapter4.verdier-localization.pf003
\begin{bukti}
	Misalkan $f: X \to Y$ suatu morfisme dalam $\mathcal{D}$. Perluas $f$
	menjadi segitiga terbedakan
	$X \xrightarrow{f} Y \to N \xrightarrow{+1}$. Ingat bahwa $N$ unik hingga
	isomorfisme. Karena itu, barisan eksak panjang suatu funktor kohomologis
	menyiratkan
% segment-id: o014.li2.chapter4.verdier-localization.q002
	\[ f \;\text{merupakan $H$-kuasi-isomorfisme} \iff \forall n, \; H(T^n N) = 0, \]
	dan ruas kanan setara dengan $N \in \Obj(\mathcal{N}_H)$. Ini membuktikan
	pernyataan.
\end{bukti}

% segment-id: o014.li2.chapter4.verdier-localization.p001
Dengan demikian, terdapat sekurang-kurangnya dua cara untuk melokalkan pada
$S_H$: menambahkan invers bagi unsur-unsur $S_H$, atau menolkan
$\mathcal{N}_H$. Ini memberikan dua sifat universal bagi lokalisasi Verdier
$Q: \mathcal{D} \to \mathcal{D}/\mathcal{N}_H
= \mathcal{D}[S_H^{-1}]$.

% segment-id: o014.li2.chapter4.verdier-localization.p002
Terakhir, kita bahas hubungan antara lokalisasi dan subkategori. Misalkan
$\mathcal{N}$ dan $\mathcal{I}$ subkategori bertriangulasi dari kategori
bertriangulasi $\mathcal{D}$, dengan $\mathcal{N}$ jenuh. Maka
$\mathcal{N} \cap \mathcal{I}$ juga merupakan subkategori bertriangulasi jenuh
dari $\mathcal{I}$. Sifat universal lokalisasi Verdier kemudian menentukan
funktor bertriangulasi
$i: \mathcal{I}/(\mathcal{N} \cap \mathcal{I}) \to \mathcal{D}/\mathcal{N}$.

% segment-id: o014.li2.chapter4.verdier-localization.p003
Perhatikan fakta sederhana berikut. Sistem
$S(\mathcal{N} \cap \mathcal{I})$ yang didefinisikan relatif terhadap
$\mathcal{N} \cap \mathcal{I} \subset \mathcal{I}$ memenuhi
% segment-id: o014.li2.chapter4.verdier-localization.q003
\begin{equation}\label{eqn:S-N-D}
	S(\mathcal{N} \cap \mathcal{I}) = S\mathcal{N} \cap \Mor(\mathcal{I}).
\end{equation}
Inklusi $\subset$ jelas. Untuk inklusi $\supset$, andaikan suatu morfisme
$f: X \to Y$ dalam $\mathcal{I}$ dapat dimasukkan ke segitiga terbedakan
$X \xrightarrow{f} Y \to N \xrightarrow{+1}$ dalam $\mathcal{D}$ dengan
$N \in \Obj(\mathcal{N})$. Korolari~\ref{prop:Cone-uniqueness} menunjukkan
bahwa $N$ isomorfik dengan suatu objek dalam $\mathcal{I}$. Karena
$\mathcal{N}$ jenuh, kita boleh mengambil
$N \in \Obj(\mathcal{N} \cap \mathcal{I})$, sehingga
$f \in S(\mathcal{N} \cap \mathcal{I})$.

% segment-id: o014.li2.chapter4.verdier-localization.pr003
\begin{proposisi}\label{prop:triangulated-i-equiv}
	Misalkan $\mathcal{N}$, $\mathcal{D}$, dan $\mathcal{I}$ seperti di atas,
	dan andaikan ``syarat resolusi'' berikut terpenuhi: untuk setiap
	$X \in \Obj(\mathcal{D})$, terdapat $U \in \Obj(\mathcal{I})$ dan suatu
	morfisme $U \to X$ yang termasuk dalam $S\mathcal{N}$. Dalam keadaan ini,
	$i: \mathcal{I}/(\mathcal{N} \cap \mathcal{I})
	\to \mathcal{D}/\mathcal{N}$ merupakan ekuivalensi kategori
	bertriangulasi.

	Kesimpulan yang sama tetap berlaku jika syarat resolusi diganti dengan:
	untuk setiap $X \in \Obj(\mathcal{D})$, terdapat
	$U \in \Obj(\mathcal{I})$ dan suatu morfisme $X \to U$ yang termasuk dalam
	$S\mathcal{N}$.
\end{proposisi}

% segment-id: o014.li2.chapter4.verdier-localization.pf004
\begin{bukti}
	Proposisi~\ref{prop:injective-localization} menyiratkan bahwa $i$ merupakan
	ekuivalensi, dan Korolari~\ref{prop:triangulated-equivalence} menunjukkan
	bahwa ekuivalensi tersebut bertriangulasi.
\end{bukti}

% segment-id: o014.li2.chapter4.verdier-localization.p004
Hasil di atas akan digunakan dalam kajian funktor turunan.
